\documentclass[../main.tex]{subfiles}

\begin{document}

\chapter*{Pendahuluan dan Orientasi Buku}
\addcontentsline{toc}{chapter}{Pendahuluan dan Orientasi Buku}

Bagian ini menjelaskan kedudukan buku ajar ini dalam mata kuliah Kriptografi,
cara memakainya, serta kerangka kurikulum yang mendasarinya. Tujuannya agar
mahasiswa dan dosen pengampu memiliki acuan yang sama sebelum masuk ke materi
bab demi bab.

\section*{1. Tujuan Buku}

Buku ajar ini disusun agar mahasiswa mampu:

\begin{enumerate}[label=\arabic*.]
    \item menjelaskan konsep kriptografi beserta urgensinya dalam dunia
          teknologi informasi;
    \item menjelaskan layanan keamanan informasi---kerahasiaan
          (\textit{confidentiality}), integritas (\textit{integrity}),
          ketersediaan (\textit{availability}), dan nirsangkal
          (\textit{non-repudiation})---serta jenis-jenis serangan terhadapnya;
    \item menggunakan landasan matematika kriptografi (aritmetika modulo,
          bilangan prima, dan invers modulo) untuk menganalisis algoritma;
    \item mengenali dan menerapkan beberapa jenis algoritma kriptografi klasik
          dan modern, baik simetris maupun asimetris, melalui contoh terhitung;
    \item menerapkan mekanisme integritas dan otentikasi---fungsi hash, MAC,
          tanda tangan digital, dan infrastruktur kunci publik---pada skenario
          nyata.
\end{enumerate}

Kelima butir di atas selaras dengan CPMK-1 (mengenal beberapa jenis algoritma
kriptografi klasik dan modern) dan CPMK-2 (memahami konsep kriptografi secara
umum dan urgensinya dalam dunia teknologi informasi).

\section*{2. Keterkaitan dengan RPS dan OBE}

Buku ini ditulis dengan kerangka \textit{Outcome-Based Education} (OBE):
capaian pembelajaran ditetapkan lebih dahulu, lalu materi, aktivitas, dan
asesmen dirancang untuk mencapainya. Setiap bab memuat komponen baku berikut.

\begin{itemize}
    \item \textbf{Tujuan Pembelajaran}---Sub-CPMK yang dicakup bab tersebut,
          dinyatakan dalam kotak \textit{Tujuan Pembelajaran}.
    \item \textbf{Materi Pokok}---uraian bagian-bagian yang menjadi isi bab,
          termasuk contoh terhitung langkah demi langkah.
    \item \textbf{Aktivitas}---langkah kerja di kelas maupun di laboratorium.
    \item \textbf{Latihan dan Refleksi}---soal yang menuntut penerapan, bukan
          sekadar hafalan.
    \item \textbf{Asesmen dan Checklist Kompetensi}---alat ukur ketercapaian
          Sub-CPMK pada akhir bab.
\end{itemize}

Pemetaan rinci antara minggu pertemuan pada RPS dan bab buku disajikan pada
bagian \textbf{Pemetaan Silabus terhadap Bab Buku}. Bobot penilaian RPS---Tugas
15\%, Kuis 10\%, Praktikum 20\%, Proyek 15\%, UTS 20\%, UAS 20\%---dijabarkan
menjadi rubrik pada Lampiran.

\section*{3. Cara Menggunakan Buku}

\begin{enumerate}[label=\arabic*.]
    \item \textbf{Baca tujuan pembelajaran lebih dahulu.} Kotak
          \textit{Tujuan Pembelajaran} di awal bab menyebut Sub-CPMK yang harus
          dicapai, sehingga Anda tahu apa yang akan diukur.
    \item \textbf{Kerjakan contoh terhitung dengan tangan.} Bagian
          \textit{Contoh Terhitung} di setiap bab memuat langkah numerik yang
          dapat Anda ulang; menuliskannya sendiri jauh lebih efektif daripada
          sekadar membaca.
    \item \textbf{Pakai aktivitas sebagai jembatan ke praktikum.} Kotak
          \textit{Aktivitas Rekomendasi} menuntun penerapan konsep bab pada
          persoalan terbuka.
    \item \textbf{Uji diri sebelum ujian.} Kerjakan \textit{Latihan dan
          Refleksi} serta \textit{Checklist Kompetensi}, lalu periksa jawaban
          terhadap kriteria rubrik pada Lampiran.
    \item \textbf{Gunakan glosarium dan daftar pustaka.} Istilah teknis
          dijelaskan pada Glosarium, sedangkan rujukan lanjutan tersedia pada
          Lampiran dan Daftar Pustaka.
\end{enumerate}

\section*{4. Ringkasan Konteks Kurikulum OBE}

Kurikulum OBE bertolak dari pertanyaan ``kompetensi apa yang harus dimiliki
lulusan?'', bukan ``materi apa yang harus diajarkan?''. Karena itu, penilaian
bukan sekadar mengukur ingatan, melainkan kemampuan menunjukkan kinerja pada
tugas yang bermakna. Sesuai prinsip tersebut, buku ini menempatkan asesmen
kompetensi, refleksi, dan rubrik pada setiap bab, serta menautkan setiap
asesmen dengan CPMK dan Sub-CPMK yang relevan. Rekapitulasi ketercapaian
kompetensi dan rekomendasi tindak lanjut disajikan pada Bab Akhir---Evaluasi,
Refleksi, dan Integrasi Kompetensi.

\end{document}
